\chapter*{Phần kết luận}
\phantomsection
\addcontentsline{toc}{chapter}{Phần kết luận}

\section*{1. Kết luận chung}
\phantomsection
\addcontentsline{toc}{section}{1. Kết luận chung}

Đề tài đã xây dựng một hệ thống thương mại điện tử thời trang tích hợp trợ lý phối đồ và thử đồ ảo ở mức sản phẩm khả dụng tối thiểu. Bên cạnh các nghiệp vụ mua sắm cơ bản, hệ thống cho phép người dùng mô tả nhu cầu bằng ngôn ngữ tự nhiên, nhận gợi ý trang phục dựa trên dữ liệu sản phẩm và tủ đồ cá nhân, đồng thời quan sát kết quả thử trang phục trên ảnh cơ thể.

\medskip

Hệ thống được tổ chức dưới dạng kho mã nguồn hợp nhất, gồm giao diện Next.js, hệ thống nghiệp vụ thương mại điện tử MedusaJS và dịch vụ AI xây dựng bằng FastAPI. Kiến trúc này phân tách rõ trách nhiệm giữa các thành phần; RabbitMQ đảm nhiệm điều phối tác vụ sinh ảnh bất đồng bộ, trong khi Server-Sent Events cập nhật trạng thái xử lý đến giao diện người dùng.

\medskip

Thực nghiệm thử đồ ảo cho thấy hai phương án xử lý đáp ứng các ưu tiên khác nhau. Với cùng bộ dữ liệu đầu vào, quy trình CatVTON hai bước mất khoảng 64 giây và duy trì khuôn mặt, tư thế cùng bối cảnh ổn định hơn. FLUX.2 hoàn thành trong khoảng 21 giây, nhanh hơn gần 67\%, tạo ảnh có chất lượng thị giác tốt nhưng làm thay đổi một số chi tiết ngoài trang phục. Kết quả này cung cấp cơ sở để lựa chọn mô hình theo yêu cầu về tốc độ hoặc mức độ bảo toàn ảnh gốc.

\section*{2. Mức độ hoàn thành}
\phantomsection
\addcontentsline{toc}{section}{2. Mức độ hoàn thành}

Các mục tiêu chính về thiết kế kiến trúc, xây dựng nghiệp vụ thương mại điện tử, triển khai AI Stylist, đồng bộ dữ liệu cho tìm kiếm ngữ nghĩa và tích hợp hai phương án thử đồ ảo đã được hoàn thành. Các luồng chức năng trọng tâm hoạt động xuyên suốt từ giao diện, hệ thống nghiệp vụ đến dịch vụ AI; kết quả thử đồ được xử lý bất đồng bộ và trả về giao diện theo thời gian thực.

\medskip

Hệ thống hiện vẫn ở phạm vi sản phẩm khả dụng tối thiểu. Những nội dung cần hoàn thiện trước khi triển khai thực tế gồm xác thực và phân quyền đầy đủ, thanh toán trực tuyến, lưu trữ ảnh dùng chung, cơ chế xử lý lại và hàng đợi lỗi, giám sát tập trung, cùng khả năng theo dõi tỷ lệ chuyển đổi từ gợi ý hoặc thử đồ sang đơn hàng.

\section*{3. Hướng phát triển}
\phantomsection
\addcontentsline{toc}{section}{3. Hướng phát triển}

Trong giai đoạn tiếp theo, ưu tiên trước hết là củng cố độ tin cậy và an toàn vận hành: chuẩn hóa cơ chế xác thực, tích hợp cổng thanh toán, chuyển ảnh sang kho đối tượng tương thích S3, bổ sung chính sách xử lý lại, hàng đợi lỗi và hệ thống giám sát tập trung. Đồng thời, luồng thanh toán cần được thống nhất với dữ liệu đơn hàng và tồn kho của MedusaJS.

\medskip

Đối với AI Stylist, cần mở rộng bộ tiêu chí đánh giá chất lượng gợi ý và ghi nhận hành vi chuyển đổi theo từng phiên làm việc. Đối với thử đồ ảo, các thực nghiệm tiếp theo nên mở rộng số lượng người mẫu, loại trang phục và điều kiện ảnh; kết quả cần được đánh giá bằng cả thời gian xử lý, chỉ số định lượng và phản hồi người dùng. Những cải tiến này sẽ giúp hệ thống tiến từ bản thử nghiệm kỹ thuật sang một sản phẩm có khả năng vận hành ổn định trong thực tế.

\clearpage
